%% ma: L10-B12-0055
%% lop: 10
%% chuong: 5
%% bai: 12
%% dang: 10.12.2
%% loai: TN4
%% muc: VD
%% dap_an: "A"
%% nguon: "chuyen-de-cac-so-dac-trung-cua-mau-so-lieu-khong-ghep-nhom-toan-10-chuong-trinh-moi.pdf (GV. Phan Nhật Linh) (P01), Bài 01 Phần I Câu 39"
%% kien_thuc: "Lớp 8 (định lí Pythagore), Lớp 10 Bài 12 (sai số, độ chính xác)"
%% trang_thai: chinh-thuc
%% ghi_chu: "Đề gốc viết 'độ dài tấm tôn được làm tròn' (thực chất là đường chéo); đã sửa lời dẫn cho rõ. Tự duyệt theo yêu cầu giáo viên 22/10/2026."
\begin{ex}
Một tấm tôn hình chữ nhật có chiều dài và chiều rộng lần lượt là $24$ m và $17$ m. Độ dài đường chéo của tấm tôn được tính rồi làm tròn đến hàng phần nghìn. Trong các số sau, số nào nhỏ nhất có thể lấy làm độ chính xác của kết quả đó?
\choice{\True $0{,}0002$}{$0{,}002$}{$0{,}005$}{$0$}
\loigiai{Đường chéo $x=\sqrt{24^2+17^2}=\sqrt{865}=29{,}41088\ldots$; làm tròn đến hàng phần nghìn được $29{,}411$.

$|x-29{,}411|\approx0{,}00012<0{,}0002$, và $d=0$ không thể (vì $x\ne29{,}411$). Chọn A.}
\end{ex}
